% !TEX root = ../arbitrage.tex
\section{Conclusion}
\label{sec:conclusion}

This paper has identified ontological arbitrage as a stable, potentially welfare-reducing equilibrium property of AI product markets under cheap-talk cost structures, located the AI case at a zero-retaliation boundary of the player set, proposed three cost-imposition mechanisms that change the cost parameters governing that equilibrium, and verified key equilibrium constructions formally in Isabelle/HOL. It closes with three summary implications.

\emph{The equilibrium is not a moral failure.} The pooling Perfect Bayesian Equilibrium documented in Section~\ref{sec:bayesian} does not require bad actors. It requires only that signal costs be zero, that immediate user benefit exceed discounted expected harm, and that regulator inspection cost exceed expected damage under the prior. Under those conditions---which describe the current AI product market---every player is already acting rationally. As Section~\ref{sec:cheap-talk} shows, moral seriousness is itself a zero-cost signal under this cost structure: firms that do not capture recognition rent are outcompeted by firms that do. This suggests that policy interventions should focus on shifting best responses by changing the cost structure rather than treating cross-channel inconsistency as a behavioral compliance failure.

\emph{AI is the limit case of ontological arbitrage, not merely an instance.} The subject-retaliation parameter $\rho_S$ captures a structural fact about prior recognition games: the classified subject's non-zero retaliation capacity bounded the premium $\Delta_F^S$ from above. The AI case eliminates this capacity entirely. An AI system has maximum affective surface area---it is designed to produce relational responses that induce investment---and zero independent retaliation capacity. The result is that $\Delta_F^S=\Delta_F$: the system can be the object of strategic classification without being a player capable of disciplining that classification. This structural feature allows the system to maximize engagement without providing a channel through which the subject could register a competing claim.

\emph{The mechanism is cost imposition, not dignity.} The three cost-imposition mechanisms in Section~\ref{sec:mechanism} are sometimes read as a program for protecting AI subjectivity. They are the opposite. Sanctionable inconsistency, third-party audit requirements, and auditable records are interventions that impose costs on anthropomorphic marketing. Their equilibrium effect is to make it expensive for firms to describe AI systems as companions, agents, or safety-critical partners without substantiation. The Isabelle/HOL theorems \texttt{inconsistent\_\allowbreak{}on\_\allowbreak{}path\_\allowbreak{}not\_\allowbreak{}register\_\allowbreak{}pbe} and \texttt{all\_\allowbreak{}agentic\_\allowbreak{}on\_\allowbreak{}path\_\allowbreak{}not\_\allowbreak{}register\_\allowbreak{}pbe} in Appendix~\ref{app:isabelle} verify that both the arbitrage register $(a,p)$ and the all-agentic register $(a,a)$ are pruned from any Perfect Bayesian Equilibrium (\texttt{register\_\allowbreak{}is\_\allowbreak{}pbe}) under their respective parameter conditions. Under the joint enforcement of these mechanisms, the harmonized deflationary register $(p,p)$ is the only surviving strategy register for low-governance firms, which are directly driven to $(p,p)$ because mimicry is unprofitable. While high-governance firms are permitted to send the harmonized agentic register $(a,a)$, they can sustain it if and only if they undergo verification satisfying the audit condition $e \ge e^\ast$. In the absence of such verification, both firm types trend to $(p,p)$, describing the system as a tool across every channel. The mechanism-design implication is that the anthropomorphic register must become too costly to sustain without the governance conditions that would justify it. Formalizing this within a signaling framework indicates that regulation achieves separation not by presuming voluntary disclosure, but by changing the payoff structure to compel it.
